\chapter{Square-Root Sigma-Point Filters}
\label{ch:bf-square-root-sigma-point}

Square-root sigma-point filters propagate a factor of the covariance rather
than the covariance matrix itself.  Their purpose is numerical: avoid avoidable
loss of symmetry and positive semi-definiteness, reduce solve instability, and
make the filtering recursion easier to monitor.  In BayesFilter, square-root
bookkeeping is a backend discipline.  It does not by itself make an
approximate likelihood exact or an autodiff gradient safe for HMC.

\section{Backend Contract}
\label{sec:bf-square-root-contract}

A square-root backend stores a factor $C_t$ such that
$P_t=C_tC_t^\top$ or an equivalent factorization.  The factor may be Cholesky,
QR-derived, SVD/eigen-derived, or another auditable representation.  The
backend contract must state:
\begin{itemize}
  \item the factor orientation and reconstruction identity;
  \item the predict update and measurement update used for the factor;
  \item whether covariance matrices are ever reconstructed;
  \item how rank loss, jitter, and PSD projection are handled;
  \item which operations define the differentiable target path.
\end{itemize}

The source SR-UKF material is useful here precisely because it exposed a
maintenance risk: a method can be called square-root while still reconstructing
full covariances and applying simplified covariance updates internally.  The
BayesFilter monograph should therefore describe the actual implemented
recursion, not only the label attached to the filter.

\section{Factor-Propagating SR-UKF Score Contract}
\label{sec:bf-srukf-factor-propagating-score-contract}

This section defines the generic square-root UKF score object used by
BayesFilter when the phrase ``factor-propagating SR-UKF analytical score'' is
used.  It is not the historical SVD sigma-point score of
Chapter~\ref{ch:bf-svd-sigma-point}.  It is also not the strict-SPD principal
matrix-square-root derivative branch in
Section~\ref{sec:bf-svd-principal-sqrt-backend}.  The backend below propagates a
declared factor and differentiates the declared factor branch.  It must not
differentiate an eigendecomposition or a principal square-root map as a hidden
substitute for the factor recursion.

\paragraph{Audit assumptions.}
\label{par:bf-srukf-audit-assumptions}
All derivatives in this section are local derivatives on one fixed branch of
the declared factor recursion.  The branch fixes the sigma-point rule, weights,
offsets, factor ranks, column order, QR pivot policy, sign convention, and the
ordered update/downdate sequence.  If any of those choices changes, the
derivative below is not reused.  For a step with state dimension \(d_x\),
observation dimension \(d_y\), sigma variable dimension \(n_a\), and factor
rank \(r_a\), the objects have shapes
\[
  \mu_a\in\mathbb R^{n_a},\quad
  C_a\in\mathbb R^{n_a\times r_a},\quad
  x^{(j)}\in\mathbb R^{d_x},\quad
  z^{(j)}\in\mathbb R^{d_y},
\]
\[
  P_{xx,\star}\in\mathbb R^{d_x\times d_x},\quad
  S_\star\in\mathbb R^{d_y\times d_y},\quad
  P_{xz,\star}\in\mathbb R^{d_x\times d_y},\quad
  K\in\mathbb R^{d_x\times d_y},\quad
  v,w,q\in\mathbb R^{d_y}.
\]
For the ordinary Gaussian innovation likelihood, the audited branch requires
\(S_\star=C_SC_S^\top\) to be symmetric positive definite on the observed
coordinates.  The implementation may evaluate \(S_\star^{-1}b\) only through
declared solves and must report solve residual and conditioning diagnostics.
The trace term in the score is square because
\(\dot S_\star^{(i)}\in\mathbb R^{d_y\times d_y}\).

\subsection{Sigma-Point Placement From A Declared Factor}
\label{sec:bf-srukf-factor-placement}

Let the sigma-point variable for one filtering step be
\[
  a_t\mid y_{1:t-1}\approx \calN(\mu_a,C_aC_a^\top),
  \qquad
  C_a\in\mathbb R^{n_a\times r_a}.
\]
The factor may be rectangular.  Rank deficiency of the represented law is
therefore allowed on a fixed-rank branch; no inverse of \(C_a\), no Cholesky
factor of \(C_aC_a^\top\), and no principal square root of \(C_aC_a^\top\) is
part of the placement contract.  For fixed standardized offsets
\(\xi^{(j)}\in\mathbb R^{r_a}\) and weights
\((w_j^{(m)},w_j^{(c)})\), the SR-UKF placement is
\begin{equation}
\label{eq:bf-srukf-factor-placement}
  a_t^{(j)}=\mu_a+C_a\xi^{(j)} .
\end{equation}
The offsets and weights define the UKF rule.  The factor \(C_a\) defines the
gauge in which that rule is placed.

For a parameter component \(\theta_i\), the first derivative of a placed point
on a fixed factor branch is
\begin{equation}
\label{eq:bf-srukf-factor-placement-first}
  \dot a_t^{(j,i)}
  =
  \dot\mu_a^{(i)}+\dot C_a^{(i)}\xi^{(j)} .
\end{equation}
The offsets are not differentiated.  A branch that changes offsets, rank,
pivoting, jitter, sign convention, or active support is a new branch and must
fail closed for analytical score admission unless that branch has its own
reviewed derivative contract.

\subsection{Moment Factors And Signed Square-Root Updates}
\label{sec:bf-srukf-moment-factor}

Let \(g_\theta\) be any smooth propagation map applied to the placed points, and
write \(u^{(j)}=g_\theta(a_t^{(j)})\).  The weighted mean and centered
deviations are
\[
  \bar u=\sum_j w_j^{(m)}u^{(j)},\qquad
  e^{(j)}=u^{(j)}-\bar u .
\]
For a covariance object that may also include an additive factor
\(G G^\top\), define positive and negative weighted stacks
\[
  E_+ =
  \begin{bmatrix}
    \{\sqrt{w_j^{(c)}}\,e^{(j)}: w_j^{(c)}>0\} & G
  \end{bmatrix},
  \qquad
  E_- =
  \begin{bmatrix}
    \{\sqrt{-w_j^{(c)}}\,e^{(j)}: w_j^{(c)}<0\}
  \end{bmatrix}.
\]
A factor-propagating SR-UKF backend chooses a declared branch
\(\mathcal B\) consisting of a positive-stack QR factorization followed by
ordered signed update/downdate primitives for the columns of \(E_-\).  The
implemented factor is
\begin{equation}
\label{eq:bf-srukf-weighted-moment-factor}
  C_u=\operatorname{SR}_{\mathcal B}(E_+,E_-),
  \qquad
  C_uC_u^\top
  =
  A_{u,\star}
  =
  E_+E_+^\top-E_-E_-^\top .
\end{equation}
The subscript \(\star\) records the implemented branch.  If a downdate fails,
a pivot changes, a sign convention changes, or a stabilizing floor is activated,
then the branch has changed.  Such an event is a diagnostic stop for the
analytical score branch unless a separate reviewed branch states the new
implemented covariance \(A_{u,\star}\) and its derivative.

The first derivative of the propagated point is
\begin{equation}
\label{eq:bf-srukf-propagated-point-first}
  \dot u^{(j,i)}
  =
  D_a g_\theta(a_t^{(j)})\,\dot a_t^{(j,i)}
  +\partial_i g_\theta(a_t^{(j)}),
\end{equation}
and the mean and covariance-object derivatives are
\begin{align}
\label{eq:bf-srukf-moment-mean-first}
  \dot{\bar u}^{(i)}
  &=
  \sum_j w_j^{(m)}\dot u^{(j,i)},\\
\label{eq:bf-srukf-moment-covariance-first}
  \dot A_{u,\star}^{(i)}
  &=
  \sum_j w_j^{(c)}
  \left(
    \dot e^{(j,i)}e^{(j)\top}
    +e^{(j)}\dot e^{(j,i)\top}
  \right)
  +\dot G^{(i)}G^\top+G\dot G^{(i)\top},
\end{align}
where \(\dot e^{(j,i)}=\dot u^{(j,i)}-\dot{\bar u}^{(i)}\).  The factor
derivative is the derivative of the same declared QR/update branch and must
reconstruct the implemented covariance derivative:
\begin{equation}
\label{eq:bf-srukf-factor-reconstruction-first}
  \dot A_{u,\star}^{(i)}
  =
  \dot C_u^{(i)}C_u^\top+C_u\dot C_u^{(i)\top}.
\end{equation}
Equation~\eqref{eq:bf-srukf-factor-reconstruction-first} is the local audit
equation for the SR-UKF factor branch.  The QR and Cholesky branch derivatives
in Chapter~\ref{ch:bf-factor-derivatives} give the primitive calculus used to
obtain \(\dot C_u^{(i)}\) when the rank, pivot, sign, and update/downdate
branch are fixed.

\subsection{Innovation Factor, Likelihood, And Score}
\label{sec:bf-srukf-innovation-score}

Apply the moment construction to the predicted state points \(x^{(j)}\) and
observation points \(z^{(j)}\).  This produces
\[
  \bar x,\qquad
  \bar z,\qquad
  P_{xx,\star},\qquad
  S_\star,\qquad
  P_{xz,\star}.
\]
The innovation factor is the declared SR factor of the observation covariance:
\begin{equation}
\label{eq:bf-srukf-innovation-factor}
  C_S C_S^\top=S_\star .
\end{equation}
On the same fixed factor branch, its first derivative reconstructs the
observation covariance derivative:
\begin{equation}
\label{eq:bf-srukf-innovation-factor-first}
  \dot S_\star^{(i)}
  =
  \dot C_S^{(i)}C_S^\top+C_S\dot C_S^{(i)\top}.
\end{equation}
For the ordinary Gaussian innovation likelihood, \(S_\star\) must be full rank
on the observed coordinates.  If it is not, the backend must either use a
separately reviewed singular-density branch or fail closed; it must not add an
unreported nugget and call the result the same likelihood.

Let \(v=y_t-\bar z\).  The solve-form likelihood uses triangular or factor
solves from the declared innovation branch:
\begin{equation}
\label{eq:bf-srukf-loglik-solve}
  C_S q=v,\qquad C_S^\top w=q,\qquad
  \ell_t
  =
  -\frac12
  \left[
    d_y\log(2\pi)+2\log|\det C_S|+q^\top q
  \right].
\end{equation}
Equivalently, \(S_\star w=v\).  The first derivative of the innovation is
\(\dot v^{(i)}=-\dot{\bar z}^{(i)}\).  Differentiating
\[
  -\frac12\{\log\det S_\star+v^\top S_\star^{-1}v\}
\]
on the same SPD solve branch uses
\(\partial_i\log\det S_\star=\tr(S_\star^{-1}\dot S_\star^{(i)})\) and
\(\partial_i S_\star^{-1}
=-S_\star^{-1}\dot S_\star^{(i)}S_\star^{-1}\).  With
\(w=S_\star^{-1}v\), the analytical score contribution is
\begin{equation}
\label{eq:bf-srukf-score-first}
  \dot\ell_t^{(i)}
  =
  -\frac12
  \left[
    \tr(S_\star^{-1}\dot S_\star^{(i)})
    +2\dot v^{(i)\top}w
    -w^\top\dot S_\star^{(i)}w
  \right].
\end{equation}
The trace term may be evaluated through the innovation factor branch, but the
object being differentiated is \(S_\star=C_SC_S^\top\) from
\eqref{eq:bf-srukf-innovation-factor}.  A derivative of a principal
matrix-square-root map is not the generic SR-UKF score route defined here.

\subsection{Filtered State And Propagated Factor Derivatives}
\label{sec:bf-srukf-filtered-state}

The UKF gain and filtered mean are
\[
  K=P_{xz,\star}S_\star^{-1},\qquad
  \mu_{t|t}=\bar x+Kv .
\]
Equivalently, \(KS_\star=P_{xz,\star}\).  Differentiating this solve equation
on the same invertible innovation branch gives the gain derivative in solve
form:
\begin{align}
\label{eq:bf-srukf-gain-first}
  \dot K^{(i)}S_\star
  &=
  \dot P_{xz,\star}^{(i)}-K\dot S_\star^{(i)},\\
\label{eq:bf-srukf-filtered-mean-first}
  \dot\mu_{t|t}^{(i)}
  &=
  \dot{\bar x}^{(i)}+\dot K^{(i)}v+K\dot v^{(i)} .
\end{align}
The filtered covariance object is
\[
  P_{t|t,\star}
  =
  P_{xx,\star}-K S_\star K^\top .
\]
The filtered square-root factor \(C_{t|t}\) is produced by the same declared
signed factor branch:
\begin{equation}
\label{eq:bf-srukf-filtered-factor}
  C_{t|t}
  =
  \operatorname{SR}_{\mathcal B_t}
  \left(C_{xx,+},C_{KS}\right),
  \qquad
  C_{t|t}C_{t|t}^\top=P_{t|t,\star},
\end{equation}
where \(C_{xx,+}\) denotes the positive state-deviation stack used for
\(P_{xx,\star}\) and \(C_{KS}\) denotes the downdated contribution whose
covariance is \(K S_\star K^\top\).  The first derivative must again reconstruct
the implemented branch:
\begin{equation}
\label{eq:bf-srukf-filtered-factor-first}
  \dot P_{t|t,\star}^{(i)}
  =
  \dot C_{t|t}^{(i)}C_{t|t}^\top
  +C_{t|t}\dot C_{t|t}^{(i)\top}.
\end{equation}
Equations~\eqref{eq:bf-srukf-filtered-mean-first} and
\eqref{eq:bf-srukf-filtered-factor-first} are the state handoff contract for a
multi-step analytical score recursion.

\subsection{Admission Boundary}
\label{sec:bf-srukf-admission-boundary}

An implementation may be called an admitted factor-propagating SR-UKF analytical
score only if its diagnostics identify the factor branch, update/downdate
events, reconstruction residuals in
\eqref{eq:bf-srukf-factor-reconstruction-first} and
\eqref{eq:bf-srukf-filtered-factor-first}, innovation solve residuals, and
score provenance.  The admitted provenance must not contain an autodiff tape
fallback, the historical SVD/eigenderivative score route, or the strict-SPD
principal-square-root derivative route.  Those routes can remain diagnostic
comparators, but they do not satisfy the generic SR-UKF score contract in this
section.

\section{HMC-Relevant Risks}
\label{sec:bf-square-root-hmc-risks}

For HMC, the relevant question is not only whether the factor recursion keeps
$P_t$ positive semi-definite in floating point.  The sampler also needs a
stable derivative of the scalar log target.  Cholesky, QR, and SVD factors have
different derivative behavior near rank loss or repeated singular values.  A
backend that is robust for value evaluation can still be fragile as a raw
tape-gradient path.

BayesFilter should therefore classify square-root sigma-point evidence in two
columns:
\begin{enumerate}[label=(\roman*)]
  \item value-recursion evidence, such as finite likelihoods, factor
    reconstruction residuals, and affine-Gaussian parity;
  \item derivative evidence, such as finite gradients, AD/custom-gradient
    parity, spectral-gap telemetry, and HMC leapfrog stability.
\end{enumerate}
Only the second column is directly relevant to HMC production readiness.

\section{When to Prefer Square-Root Forms}
\label{sec:bf-square-root-choice}

Square-root forms are natural when observation noise is small, latent
covariances are nearly singular, or long scans accumulate asymmetry.  They are
also useful for large LGSSMs where factor diagnostics expose numerical problems
earlier than a final log likelihood does.  They should not be treated as a
universal answer to nonlinear filtering: if the moment closure is poor, the
filter is still approximating the wrong target; if the derivative of the
factorization is unstable, HMC can still fail.

\section{CUT With Square-Root Backends}
\label{sec:bf-cut-square-root}

The conjugate unscented transform in Section~\ref{sec:bf-cut-rules} is a
point-and-weight rule.  It can be combined with any square-root factor backend
that maps standardized points $\xi^{(j)}$ into state points
\[
  \chi^{(j)}=m+C\xi^{(j)},\qquad CC^\top=P_\star .
\]
The experimental
\texttt{/home/chakwong/python/src/dsge\_hmc/filters/CUTSRUKF.py} module uses
positive CUT4-G weights together with QR square-root covariance updates.  It also
keeps rank-deficient process noise as a rectangular factor $G$ with
$Q=GG^\top$, avoiding an invalid Cholesky factorization of a singular $Q$.

An SVD/eigen factor can be used in the same location:
\[
  P=U\Lambda U^\top,\qquad C=U\Lambda_+^{1/2}.
\]
This may make value evaluation more robust when $P$ is singular or nearly
singular.  It does not change CUT's quadrature degree; it changes the numerical
factor used to place the CUT points.  Therefore a stable CUT-SVD filter must
report the same factor diagnostics as an SVD sigma-point filter: the covariance
actually reconstructed, active floors or rank truncations, spectral gaps, sign
or order choices, and fallback branches.  In short, CUT can be combined with
SVD, but the result is a higher-order quadrature rule on top of an SVD factor
policy, not a new theorem that removes spectral derivative risk.
